\documentclass[11pt, reqno]{amsart}

% --- Fonts, encoding, geometry ---
\usepackage[utf8]{inputenc}
\usepackage[T1]{fontenc}
\usepackage[margin=1in]{geometry}

% --- AMS math stack ---
\usepackage{amsmath, amssymb, amsthm, mathtools}

% --- Graphics, tables, hyperlinks ---
\usepackage{graphicx}
\usepackage{booktabs}
\usepackage[colorlinks=true, linkcolor=blue!60!black, citecolor=blue!60!black, urlcolor=blue!60!black]{hyperref}
\usepackage[capitalise, noabbrev]{cleveref}

% --- Bibliography ---
\usepackage[numbers, sort&compress]{natbib}

% --- Theorem environments (shared counter, sectioned) ---
\theoremstyle{plain}
\newtheorem{theorem}{Theorem}[section]
\newtheorem{lemma}[theorem]{Lemma}
\newtheorem{proposition}[theorem]{Proposition}
\newtheorem{corollary}[theorem]{Corollary}

\theoremstyle{definition}
\newtheorem{definition}[theorem]{Definition}
\newtheorem{example}[theorem]{Example}
\newtheorem{conjecture}[theorem]{Conjecture}

\theoremstyle{remark}
\newtheorem{remark}[theorem]{Remark}

% --- Custom operators (add as needed) ---
\DeclareMathOperator{\Aut}{Aut}
\DeclareMathOperator{\End}{End}
\DeclareMathOperator{\Hom}{Hom}
\DeclareMathOperator{\id}{id}
\DeclareMathOperator{\im}{im}
\DeclareMathOperator{\rank}{rank}
\DeclareMathOperator{\Tr}{Tr}

% --- Number sets ---
\newcommand{\N}{\mathbb{N}}
\newcommand{\Z}{\mathbb{Z}}
\newcommand{\Q}{\mathbb{Q}}
\newcommand{\R}{\mathbb{R}}
\newcommand{\C}{\mathbb{C}}
\newcommand{\F}{\mathbb{F}}

% --- Metadata ---
\title{Title}
\author{Author Name}
\address{Institution, Department}
\email{author@institution.edu}
\date{\today}

\subjclass[2020]{Primary ; Secondary }
\keywords{keyword one, keyword two}

\begin{document}

\begin{abstract}
Abstract text. State the main result, the setting, and the key method in one paragraph.
\end{abstract}

\maketitle

\section{Introduction}\label{sec:intro}

Motivate the problem, state the main theorem, sketch the proof strategy, and
outline the paper.

\begin{theorem}[Main]\label{thm:main}
Statement of the main theorem.
\end{theorem}

\section{Preliminaries}\label{sec:prelim}

Definitions, notation, cited background results.

\begin{definition}\label{def:widget}
A \emph{widget} is \ldots
\end{definition}

\section{Auxiliary Results}\label{sec:aux}

Lemmas used in the main proof.

\begin{lemma}\label{lem:key}
Statement of the key lemma.
\end{lemma}
\begin{proof}
Proof body. Each step cites a theorem, lemma, definition, or named inference
rule (see \cref{def:widget}).
\end{proof}

\section{Proof of the Main Theorem}\label{sec:proof}

\begin{proof}[Proof of \cref{thm:main}]
Proof body, structured as stated in \cref{sec:intro}. Cite auxiliary results
with \verb|\cref{lem:key}|.
\end{proof}

\section{Open Questions}\label{sec:open}

\begin{conjecture}\label{conj:one}
A conjecture the paper leaves open.
\end{conjecture}

\bibliographystyle{amsplain}
\bibliography{references}

\end{document}
